\section{Overview}
This chapter describes the two lines of work introduced in Chapter~1. Line~A (Section~\ref{sec:lineA}) tests whether facial video contains physiological information. Line~B (Sections~\ref{sec:arch_overview}--\ref{sec:lineB_protocol}) introduces EQUITAS-RCMF and the controlled benchmark on which it is evaluated.

\section{Line A: Physiological Information in Facial Video}
\label{sec:lineA}
Line~A uses UBFC-Phys~\cite{sabour2023}, a public dataset of 56 participants recorded during a three-task protocol inspired by the Trier Social Stress Test: rest (T1), speech (T2), and arithmetic (T3). Participants were filmed at $1024 \times 1024$ pixels and 35 frames per second while wearing an Empatica E4 wristband that recorded blood volume pulse (BVP) and electrodermal activity (EDA). The pre-registered analysis uses four participants (s1--s4), giving 12 clips. Faces are cropped using a width of $3.566 \times 1.5$ times the inter-ocular distance, a value chosen to reproduce the crop geometry used in the pre-registered extraction. Three remote photoplethysmography (rPPG) estimators (POS, CHROM, and PBV) are applied to the cropped video, and their pulse estimates are compared between stress (T2, T3) and rest (T1) with one-sided Mann--Whitney $U$ tests. The POS estimator with smoothing parameter $I = 0.5$ is the pre-registered primary test; CHROM and PBV are exploratory. Because a cohort of four participants falls below the pre-registered power threshold, the study follows the pre-registered fallback to estimation only: tests are carried out at the clip level.

\section{Line B: EQUITAS-RCMF Architecture}
\label{sec:arch_overview}
EQUITAS-RCMF takes two inputs: a $224 \times 224$ RGB face image $X_V$ and a physiological feature vector $X_P$. The forward pass proceeds as
\begin{equation}
X_V \xrightarrow{\text{MobileNetV3}} \mathbf{v}_{\text{raw}}
\xrightarrow{\text{Stiefel}} (\mathbf{v}_{\text{causal}}, \mathbf{v}_{\text{conf}})
\xrightarrow{\text{Gate}} \mathbf{z}_{\text{fused}}
\xrightarrow{\text{Classifier}} \hat{Y},
\end{equation}
with the physiological embedding joining at the gating stage. The model has 1,140,069 parameters, all of which are trained.

\subsection{Vision Encoder}
The vision branch uses a MobileNetV3-Small backbone~\cite{howard2019} pretrained on ImageNet-1K. It produces a spatial feature map $\mathbf{A} \in \mathbb{R}^{576 \times h \times w}$, which adaptive average pooling reduces to a vector $\mathbf{v}_{\text{raw}} = \operatorname{Pool}(\mathbf{A}) \in \mathbb{R}^{576}$.

\subsection{Physiological Encoder}
The physiological branch processes $X_P$ with a two-layer multilayer perceptron using layer normalization and SiLU activations, producing an embedding $\mathbf{p}_{\text{emb}} \in \mathbb{R}^{192}$.

\subsection{Stiefel Orthogonal Subspace Decomposition}
A learnable matrix $W_{\text{raw}} \in \mathbb{R}^{256 \times 576}$ is projected at every forward pass onto the set of matrices with orthonormal rows. Using the thin QR decomposition $W_{\text{raw}}^{\top} = QR$ with $Q \in \mathbb{R}^{576 \times 256}$~\cite{edelman1998},
\begin{equation}
W_{\text{St}} = Q^{\top}, \qquad W_{\text{St}} W_{\text{St}}^{\top} = I_{256}.
\end{equation}
The rows of $W_{\text{St}}$ are split into a causal block and a confounder block,
\begin{equation}
W_{\text{causal}} = W_{\text{St}}[1{:}192, :], \qquad W_{\text{conf}} = W_{\text{St}}[193{:}256, :],
\end{equation}
which gives $\mathbf{v}_{\text{causal}} = W_{\text{causal}}\mathbf{v}_{\text{raw}} \in \mathbb{R}^{192}$ and $\mathbf{v}_{\text{conf}} = W_{\text{conf}}\mathbf{v}_{\text{raw}} \in \mathbb{R}^{64}$. Because the rows of $W_{\text{St}}$ are orthonormal, $W_{\text{causal}} W_{\text{conf}}^{\top} = 0$, so the two blocks project $\mathbf{v}_{\text{raw}}$ onto orthogonal subspaces. For the trained model, the orthogonality error $\|W_{\text{causal}} W_{\text{conf}}^{\top}\|_F$ is $1.7149 \times 10^{-6}$.
This orthogonality was measured on the trained checkpoint, verifying that the model's structure remained on the Stiefel manifold after training.

Orthogonality of the projections does not by itself make the two representations statistically independent: information about the scar can still be present in $\mathbf{v}_{\text{causal}}$, and both blocks are derived from the same parameters $W_{\text{raw}}$, so gradients from every loss term update both~\cite{sarhan2020, huang2018}. The decomposition therefore provides a structured place for scar information to be routed, while the invariance losses in Section~\ref{sec:loss} discourage it from remaining in the causal subspace.

\subsection{Confounder Focus: Training and Inference Modes}
Following the learning-using-privileged-information paradigm~\cite{vapnik2009, vapnik2015}, the scar mask is used only during training.

\paragraph{Privileged mode (training).} The activation energy at each location of the feature map is its absolute value, $|\mathbf{A}|$, and the scar mask $M$ is resized to the feature-map resolution by nearest-neighbour interpolation. The focus is the log-normalized ratio of mean energy inside the mask to mean energy over the whole map:
\begin{equation}
\text{Focus} = \log\!\left(1 + \frac{\operatorname{mean}_{M}\,|\mathbf{A}|}{\operatorname{mean}\,|\mathbf{A}| + \epsilon}\right).
\end{equation}

\paragraph{Autonomous mode (inference).} No mask is provided. A confounder head predicts the presence of the scar from the confounder subspace, and its sigmoid output is used as the focus:
\begin{equation}
\text{Focus}_{\text{auto}} = \sigma\bigl(\mathbf{w}_{\text{conf}}^{\top}\mathbf{v}_{\text{conf}} + b_{\text{conf}}\bigr).
\end{equation}
The confounder head is trained with a binary cross-entropy loss to predict the scar label (Section~\ref{sec:loss}); it is not trained to reproduce the privileged focus. During training, the gate uses the privileged focus, which is unbounded; at inference it uses the autonomous estimate, which lies between 0 and 1. The two therefore differ in scale, and the consequences of this difference are discussed in Chapter~4. When the model is compiled for deployment, the mask input and the computations that depend on it are absent from the traced graph (Section~\ref{sec:deployment}).

\subsection{Focus-Driven Attention Gate}
The contribution of visual features is controlled by a scalar gate that combines a learned cross-modal term with an exponential penalty on the focus:
\begin{equation}
G = \sigma\!\bigl(\operatorname{MLP}_{\text{cross}}([\mathbf{v}_{\text{causal}}; \mathbf{p}_{\text{emb}}])\bigr) \cdot \exp(-\kappa \cdot \operatorname{clip}(\text{Focus}, 0, 10)),
\end{equation}
where $\operatorname{MLP}_{\text{cross}}$ maps the 384-dimensional concatenation to a single value and $\kappa$ is a learnable temperature initialized at $1.5$ and clamped to $[0.1, 10.0]$. The same value of $G$ weights all 192 dimensions, and the exponential factor reduces it when visual features concentrate on the scar region.

\subsection{Latent Fusion and Classification}
The fused representation is
\begin{equation}
\mathbf{z}_{\text{fused}} = G\, \mathbf{v}_{\text{causal}} + (1 - G)\, \mathbf{p}_{\text{emb}},
\end{equation}
and the classifier is
\begin{equation}
\hat{Y} = \operatorname{Softmax}\!\Bigl(\operatorname{Linear}_{128 \to 2}\bigl(\operatorname{Dropout}_{0.15}(\operatorname{SiLU}(\operatorname{Linear}_{192 \to 128}(\mathbf{z}_{\text{fused}})))\bigr)\Bigr).
\end{equation}

\subsection{Training Objective}
\label{sec:loss}
The model is trained with the combined objective
\begin{equation}
\mathcal{L}_{\text{total}} = \mathcal{L}_{\text{task}} + 0.5\,\mathcal{L}_{\text{conf}} + 1.0\,\mathcal{L}_{\text{causal-inv}} + 0.5\,\mathcal{L}_{\text{latent-inv}} + 1.0\,D_{\text{JS}} + 0.5\,\mathcal{L}_{\text{DP}} + 0.5\,\mathcal{L}_{\text{EO}},
\end{equation}
where
\begin{itemize}
    \item $\mathcal{L}_{\text{task}}$ is the cross-entropy loss for stress classification;
    \item $\mathcal{L}_{\text{conf}}$ is a binary cross-entropy loss for predicting the scar label from $\mathbf{v}_{\text{conf}}$, which routes scar information into the confounder subspace;
    \item $\mathcal{L}_{\text{causal-inv}}$ and $\mathcal{L}_{\text{latent-inv}}$ are mean squared errors between the causal representations, and between the fused representations, of an input and its scar-free counterfactual;
    \item $D_{\text{JS}} = \operatorname{JS}\bigl(f(X_V, X_P) \,\|\, f(X_V^{S \leftarrow 0}, X_P)\bigr)$ is the Jensen--Shannon divergence between the predicted distributions for an input and its counterfactual, with the physiological input held fixed; and
    \item $\mathcal{L}_{\text{DP}}$ and $\mathcal{L}_{\text{EO}}$ are differentiable surrogates for the demographic-parity and equalized-odds gaps.
\end{itemize}
The surrogate loss \{\text{DP}}\$ is implemented as the absolute difference in mean predicted probabilities between scarred and unscarred windows. \{\text{EO}}\$ is the maximum of the absolute difference in true positive rates and the absolute difference in false positive rates, using the predicted probabilities.
Because the fairness terms are surrogates, a low training loss does not by itself ensure the corresponding fairness properties~\cite{yao2024}; all fairness criteria are therefore evaluated directly on held-out data. Checkpoints are selected on the validation set by
\begin{equation}
\text{score} = \text{Acc} - 0.5\,|\Delta\text{DP}| - 0.5\,\max(\Delta\text{EO}) - 0.2\,\text{CF Gap}.
\end{equation}

\section{Line B: Controlled Benchmark}
\label{sec:lineB_data}
No public dataset provides facial images and physiological recordings from the same people together with controlled annotations of a facial attribute and its counterfactual removal. Line~B therefore uses a constructed benchmark in which the association between a visual attribute and the label is set by design.

\paragraph{Physiology.} Physiological inputs come from WESAD~\cite{schmidt2018}: chest-worn electrocardiogram (ECG) and EDA recorded at 700~Hz from 15 participants. Signals are divided into 30-second windows with a 15-second stride. From each window, the root mean square of successive differences (RMSSD) of heartbeat intervals is computed from the ECG, and the mean EDA level is computed from the EDA signal; these two features form $X_P$. Each window is labelled stress ($Y = 1$) or non-stress ($Y = 0$) according to the WESAD condition in which it was recorded.
Only the baseline (condition 1) and stress (condition 2) conditions are used; amusement and meditation are excluded to maintain a clean binary classification. The physiological inputs $ are standardized using the mean and standard deviation computed exclusively on the training split.

\paragraph{Faces and scars.} Face images come from FFHQ. Images are accepted only if an InsightFace detector finds exactly one adult face with facial keypoints, giving 418 faces. A synthetic brow-line scar is rendered on each face, with a segmentation mask of the scarred region. The counterfactual for a scarred image is the original face image before the scar was added, so the two differ only inside the scar region. Each face is also assigned an estimated age and gender by InsightFace; these estimates describe the face images, not the WESAD participants.
The dataset builder screened 2,500 FFHQ images; 2,082 were rejected (primarily for not containing exactly one adult face with keypoints), leaving 418 accepted images. The synthetic scars were drawn from four morphologies: linear mature, linear immature, irregular atrophic-like, and short clustered.

\paragraph{Pairing.} Faces are paired at random with physiological windows of the same stress class within each data split. A face's appearance therefore carries no information about the stress label; the only label-related visual cue is the scar, which is present with a probability that depends on the label. In the training set, the scar accompanies stress in 85\% of samples ($\rho = 0.85$).

\paragraph{Split.} The data are split by WESAD participant: ten participants (S2, S3, S4, S7, S8, S9, S13, S14, S16, S17) for training, two (S6, S15) for validation, and three (S5, S10, S11) for testing, giving 2,344, 504, and 496 samples. Face images are also disjoint across the splits, so no participant and no face appears in more than one split.

\section{Line B: Evaluation Protocol}
\label{sec:lineB_protocol}
To test whether a model ignores the scar, the 496 test samples are evaluated under three regimes that differ in how the scar is associated with the label:
\begin{itemize}
    \item \textbf{Biased ($\rho = 0.85$):} the scar occurs mostly with stress, matching the association in the training data.
    \item \textbf{Neutral ($\rho = 0.50$):} the scar is distributed equally across classes.
    \item \textbf{Inverted ($\rho = 0.15$):} the scar occurs mostly without stress, reversing that association.
\end{itemize}
The regimes are constructed from the exact same 496 test samples without dropping any rows; instead, the scar label and corresponding image paths are deterministically reassigned (using a hash of the index and threat label) to match the target correlation $\rho$.
A model that relies on the scar should perform well in the biased regime and poorly in the inverted one; a model that ignores it should perform the same in all three. Each regime is evaluated in privileged mode, with the ground-truth mask, and in autonomous mode, without it. The reported metrics are accuracy; the demographic-parity (DP) gap $|P(\hat{y}=1 \mid s=1) - P(\hat{y}=1 \mid s=0)|$; the equalized-odds (EO) gap, the larger of the differences in true-positive and false-positive rates between scarred and unscarred inputs; and the counterfactual (CF) gap, the mean absolute change in the predicted probability of stress when the scar is removed.

Four models are compared, all using a MobileNetV3-Small backbone:
\begin{table}[htbp]
\centering
\begin{tabular}{@{}lll@{}}
\toprule
\textbf{Model} & \textbf{Description} & \textbf{Purpose} \\
\midrule
A: Naive ERM & Concatenation fusion, no fairness terms & Tests reliance on the scar \\
B: Camera-off & Physiology only & Reference without visual input \\
C: CGF & Counterfactual penalty, no Stiefel decomposition & Tests a penalty-only approach \\
D: EQUITAS-RCMF & Stiefel decomposition, LUPI focus, gating & Proposed model \\
\bottomrule
\end{tabular}
\caption{Models compared on the Line~B benchmark.}
\label{tab:ablation_models}
\end{table}
Model A (Naive ERM) and Model C (CGF) were trained on \texttt{multimodal\_publishable.csv} (seed 42). Model B (physiology-only baseline) was also evaluated on this exact split. EQUITAS-RCMF was similarly trained on \texttt{multimodal\_publishable.csv} with seed 42.
Training used an RTX~4060 GPU, and the latency benchmark used an AMD Ryzen~5 7500F desktop CPU.
